\section{连续性与整体保证}
\label{sec:02-Calculus/02-Limits-and-Continuity/04}

\subsection{没有公式时，凭什么说方程有解}

考虑一个简单的三次方程：

\[
x^3+x-1=0.
\]

没有求根公式能漂亮地给出精确解——三次方程的卡尔丹公式的确存在，但在这里用起来远不如直接思考。

换个思路：令 \(p(x)=x^3+x-1\)。算两个点：

\[
p(0)=0^3+0-1=-1,\qquad p(1)=1^3+1-1=1.
\]

一个负，一个正。如果 \(p\) 的图像是一条不间断的曲线，那么从 \(-1\) 走到 \(1\) 的路上，必定穿过横轴。也就是说，\((0,1)\) 之间至少有一个根。

这个论证用到了一个你甚至没有仔细想过的假设：函数的图像不会 ``断开''。多项式 \(p(x)=x^3+x-1\) 是连续的，图像确实不间断，所以论证成立。

但换个情况呢？假设某函数 \(f\) 满足 \(f(0)=-1\)，\(f(1)=1\)，但在 \(x=0.5\) 处函数值突然跳到 10 然后又跳回来——那么 \(f\) 的图像完全可以从 \(-1\) 出发，跳到高空中，再从另一边落回 \(1\)，路上从未经过 0。``两端异号则中间必有零点''这个直觉，需要函数的连续性质作为担保。

\subsection{连续：点值与附近趋势接上了}

函数 \(f\) 在点 \(a\) 连续，若

\[
\lim_{x\to a}f(x)=f(a).
\]

这一行包含三件事：\(f(a)\) 有定义，极限存在，二者相等。缺任何一项都不连续。

举个例子，函数

\[
f(x)=\frac{x^2-1}{x-1}
\]

在 \(x=1\) 处没有定义，所以谈不上连续——尽管 \(\lim_{x\to1}f(x)=2\) 存在。这属于可去间断：补上 \(f(1)=2\) 就能修复。若补成其他值，极限虽存在却与点值不相等，仍然是间断。

再看 \(1/x\) 在 \(x=0\) 附近：它根本无界，极限不存在且不是有限的。\(\sin(1/x)\) 则在 \(x=0\) 附近无限振荡，极限也不存在。阶跃函数——比如 \(x<0\) 时为 \(-1\)，\(x\ge0\) 时为 \(1\)——在 0 处有跳跃间断：左右极限都存在，但不相等。

把所有不连续都叫 ``断了一下'' 会掩盖根本不同的失败机制：有的时候是没定义，有的时候是极限不存在，有的是极限与点值不匹配。知道哪种失败，才知道有没有办法修复。

连续性也可从误差控制的角度看：输入足够接近 \(a\)，输出就足够接近 \(f(a)\)。与一般极限相比，这里不再排除点值，而且目标值由函数本身确定——不需要额外指定一个 \(L\)。

\subsection{连续函数可以安全地层层代入}

多项式处处连续；有理函数在分母非零处连续；指数、对数、三角函数在各自定义域内连续。连续函数的和、积、合法商与复合仍连续——这些性质来自极限运算法则的直接搬运。

因此求

\[
\lim_{x\to0}e^{\sin(x^2)}
\]

可以逐层代入：\(x^2\to0\)，\(\sin\) 连续给 \(\sin(x^2)\to\sin0=0\)，\(\exp\) 连续给最终极限 \(e^0=1\)。这里的 ``直接代入'' 是复合连续性的压缩应用——每一步都因为对应函数连续而合法。

\subsection{介值定理：局部无跳跃变成整体穿越}

把刚才三次方程的直觉写成严谨陈述，就是介值定理（Intermediate Value Theorem）：

若 \(f\) 在闭区间 \([a,b]\) 连续，且数 \(N\) 位于 \(f(a)\) 与 \(f(b)\) 之间，那么存在 \(c\in[a,b]\) 使 \(f(c)=N\)。

对于 \(p(x)=x^3+x-1\)，\(p(0)=-1\)，\(p(1)=1\)，取 \(N=0\)，定理断言 \((0,1)\) 内至少有一个根。注意定理给的是存在性，不直接给公式。它也没有保证唯一性——在本例中，因为 \(p'(x)=3x^2+1>0\)，函数严格递增，才能进一步得到唯一根。

阀门的类比更直观。水箱液位 \(h(t)\) 如果连续，8 点时低于警戒线，9 点高于警戒线，那么这一个小时内必有某一时刻恰好到达警戒线。定理保证的是这个时刻存在，不意味着每隔五分钟采样的记录中恰好出现了该数值；若传感器读数经过取整处理，读数还可能直接跳过警戒值。应用介值定理时，连续的是物理状态、测量信号还是离散记录，必须分清。

连续条件不能删。阶跃函数可以从负值跳到正值而从未取 0。定义域的连通性也重要——若定义域只含两个孤立的端点，端点值跨过某数并不迫使中间存在输入。

\subsection{极值定理：闭区间上连续函数一定取到最大最小值}

另一个全局保证是极值定理（Extreme Value Theorem）：

若 \(f\) 在闭有界区间 \([a,b]\) 上连续，则它一定取得最大值和最小值。

这里的 ``取得'' 比 ``有上界'' 更强：确实存在点 \(x_{\max},x_{\min}\) 达到极值。

每个条件都不可或缺：

\begin{itemize}
\tightlist
\item
  \(f(x)=x\) 在开区间 \((0,1)\) 上有上确界 1，却没有最大值——因为 1 不在区间内，最靠近 1 的点永远可以 ``再靠近一点点''；
\item
  \(f(x)=1/x\) 在 \((0,1]\) 上连续，却无上界——当 \(x\to0^+\) 时函数值冲向无穷；该区间不是闭的；
\item
  在闭区间上故意挖掉一个点（在某个 \(x_0\) 处重新定义函数值使之上跳），也可能让上确界不被取到——不连续处可以 ``跳过'' 本该取到的值。
\end{itemize}

更高层的表述是：连续函数把 ``有界且包含所有边界点'' 的集合映成同样性质的集合，这类集合在数学上叫紧致集，闭区间只是实数中最常见的紧致集。

\subsection{一点连续不等于全局温和}

函数在每一点连续，不表示全域有界或变化缓慢。\(e^x\) 在实数上处处连续却无界——随着 \(x\to\infty\) 它一路飙升。\(\sin(x^2)\) 处处连续却在远处振荡越来越快，导数无界。

原因在于：单点连续只保证在足够小的邻域内输出波动可控。``足够小'' 这个尺度可以逐点变化——在 \(e^x\) 的远处，你需要把输入框得很窄才能把输出波动压到给定范围。没有紧致性、单调性或其他全局结构的加持，逐点连续性本身不承诺全局的温顺。

反过来看介值和极值定理的意义：它们正是把 ``逐点'' 的连续性，与 ``闭区间'' 的全局结构结合后，产出的整体结论。连续性是每个点的性质，定理是把局部性质编织成全局保证的那根针。

延伸阅读：MIT OpenCourseWare 18.01SC: Limits and Continuity。
